%%%%%%%%%%%%%%%%%%%%%%%%%%%%%%%%%%%%%%%%%%%%%%%%%%%%%%%%%%%%%%%

%%Modelo desenvolvido pelo Professor JR Ferreira do ICET - UFAM.

%Data: 28/10/2023.

%%%%%%%%%%%%%%%%%%%%%%%%%%%%%%%%%%%%%%%%%%%%%%%%%%%%%%%%%%%%%%%

\documentclass[a4paper,11pt]{article} %Selecionando a classe que gera artigos.
\usepackage{background}
\backgroundsetup{contents={\includegraphics[scale=1]{figures/cabecalho.jpg}},scale=1,placement=top,opacity=1,position={7.975cm,1.35cm}, vshift=45}

\input{preambulo} %Pacotes principais

%%%%%%%%%%%%%%%%%%%%%%%%%%%%%%PACOTES SECUNDÁRIOS%%%%%%%%%%%%%%%%%%%%%%%%%%%%%%



\pagestyle{empty}%retirando numeração
\usepackage{fancyhdr}
\pagestyle{fancy}
\renewcommand{\headrulewidth}{0pt}
\lhead{}
\chead{}
\rhead{}
\lfoot{}
\cfoot{\includegraphics[scale=.35]{figures/re.png}}
\rfoot{}




\begin{document} %Begin Inicia o Documento



\input{cabecalho} %Alterar Título, autores, e apagar tabela !

%%%%%%%%%%%%%%%%%%%%%%%%%%%%%%TEXTOS%%%%%%%%%%%%%%%%%%%%%%%%%%%%%%

\begin{abstract}
\noindent A primeira página deve constar apenas: Título, Autores, Resumo e Palavras-Chave. O Resumo é elemento obrigatório, não devendo ultrapassar 200 (duzentas) palavras. No Resumo deve ser apresentado o problema em estudo, os métodos utilizados, os resultados mais importantes e as conclusões obtidas. O Resumo não deve extrapolar para a segunda página. Xxxx xxxxx xxxxxxxx xxxxxxxxxxx xxx xxxxxxxx xxxxxxx xxxxx xxxxx xxxx xxx xxxxx xxxxxxxx xxxxxxxxxxx xxx xxxxxxxx xxxxxxx xxxxx xxxxx xxxx xxx xxxxx xxxxxxxx xxxxxxxxxxx xxx xxxxxxxx xxxxxxx xxxxx xxxxx xxxx xxx xxxxx xxxxxxxx xxxxxxxxxxx xxx xxxxxxxx xxxxxxx xxxxx xxxxx xxxx xxx xxxxx xxxxxxxx xxxxxxxxxxx xxx xxxxxxxx xxxxxxx xxxxx xxxxx xxxx xxx xxxxx xxxxxxxx xxxxxxxxxxx xxx xxxxxxxx xxxxxxx xxxxx xxxxx xxxx xxx xxxxx xxxxxxxx xxxxxxxx xxx xxxxxxxx xxxxxxx xxxxx xxxxx xxxx xxx xxxxx xxxxxxxx xxxxxxxxxxx xxx xxxxxxxx xxxxxxx xxxxx xxxxx xxxx xxx xxxxx xxxxxxxx xxxxxxxxxxx xxx xxxxxxxx xxxxxxx xxxxx xxxxx xxxx xxx xxxxx xxxxxxxx xxxxxxxxxxx xxx xxxxxxxx xxxxxxx xxxxx xxxxx xxxx xxxx xxx xxxxx xxxxxx xxxxxx xxx xxxxx xxxxxxx xxxxx xxxxx xxxx xxx xxxxx xxxxxxxx xxxxxxxxxxx xxx xxxxxxxx xxxxxxx xxxxx xxxxx xxxx. Xxxx xxxxx xxxxxxxx xxxxxxxxxxx xxx xxxxxxxx xxxxxxx xxxxx xxxxx xxxx xxx xxxxx xxxxxxxx xxxxxxxxxxx xxx xxxxxxxx xxxxxxx xxxxx xxxxx xxxx xxx xxxxxxxx xxxxxxx xxxxx xxxx xxxx xxxxx xxxx.
\end{abstract}


\vspace{0.5cm}
\noindent\textbf{Palavras-Chave:}: Três a cinco Palavras-Chave separadas por ponto e vírgula (;), procurando não repetir palavras do título, escritas em letras minúsculas. (Times New Roman, 11, Justificado).


\section{INTRODUÇÃO}

O Resumo Expandido do XXXII CONIC poderá conter até 4 páginas, incluindo as referências, sendo a primeira destinada ao Título, Autores, Resumo e Palavras-Chave. Para padronização, os autores deverão \textcolor{red}{utilizar este arquivo} devidamente formatado. A fonte a ser usada em todo o trabalho deverá ser “calibri”. O corpo do texto deverá ser digitado usando fonte calibri de tamanho 12, incluindo o Título (tamanho 14) e Subtítulos (tamanho 12), espaçamento 1,0 entre linhas e alinhamento justificado.

Padronizar o formato utilizado nos trabalhos é essencial para a correta edição dos Anais do Evento, bem como a disponibilização e acesso ao Resumo Expandido através do DOI. Espera-se que, ao utilizar os estilos pré-definidos deste documento, a formatação de seu Resumo Expandido seja facilitada. Neste sentido, observe as instruções e formate o resumo de acordo com os padrões definido. 

O trabalho que não for submetido de acordo com este modelo não seguirá para o processo de avaliação.

Lembre-se que a introdução de um Resumo Expandido desempenha um papel fundamental ao estabelecer o contexto, apresentar o problema de pesquisa e fornecer uma visão geral do estudo. A introdução deve ser escrita de forma clara e envolvente, despertando o interesse do leitor e fornecendo informações essenciais para contextualizar o restante do resumo expandido. Certifique-se de que a introdução seja concisa e direta, evitando detalhes excessivos ou informações irrelevantes. Ela serve como uma base sólida para o leitor entender o propósito e a relevância da pesquisa que será apresentada.

%No máximo 20 linhas, evitar divagações, utilizando bibliografia apropriada para formular os problemas abordados e a justificativa da importância do assunto, deixando claro a(s) hipótese(s) e o(s) objetivo(s) do trabalho. (Times New Roman, 11, Justificado, Máximo 20 linhas). De acordo com \cite{alves} podemos...

\section{OBJETIVO GERAL}

O objetivo geral em um Resumo Expandido é uma declaração concisa que descreve o propósito principal do estudo ou pesquisa em questão. Ele deve fornecer uma visão geral clara e abrangente do que o estudo visa alcançar. 

\section{METODOLOGIA}

A seção de metodologia em um Resumo Expandido deve fornecer uma descrição concisa e clara dos métodos e procedimentos utilizados para conduzir a pesquisa. Esta seção é fundamental, pois permite que os leitores compreendam como o estudo foi realizado e avaliem a validade e a robustez dos resultados apresentados. A seção de metodologia deve ser clara e detalhada o suficiente para que outros pesquisadores possam entender como o estudo foi conduzido e possam replicá-lo, se necessário. No entanto, lembre-se de que um Resumo Expandido tem limitações de espaço. Portanto, a descrição da metodologia deve ser sucinta e focada nos aspectos mais relevantes para o entendimento da pesquisa.

\section{RESULTADOS E DISCUSSÃO}

A seção de resultados e discussão em um Resumo Expandido deve fornecer uma síntese clara e concisa dos principais achados da pesquisa, bem como uma análise e interpretação desses resultados. Esta seção é essencial para comunicar as descobertas do estudo aos leitores. Lembre-se de que a seção de resultados e discussão deve ser clara e lógica, seguindo uma progressão natural dos resultados para a discussão. Concentre-se nos resultados mais importantes e na interpretação crítica deles. Mantenha o foco no objetivo da pesquisa e nas questões de pesquisa levantadas. Além disso, evite especulações não fundamentadas e baseie suas análises em evidências sólidas.

\begin{figure}[H]
	\centering
	\includegraphics[width=10cm,angle=0]{figures/pizza.jpg}
	\caption{Descrição da figura. Utilizar fonte Calibri tamanho 11. (se houver).}
	\label{fig:pizza}
\end{figure}

\begin{table}[H]
	\centering
	\caption{Título da Tabela. Utilizar fonte Calibri tamanho 11. (se houver)}
	\label{tab:tabela1}
	\begin{tabular}{|l|l|l|l|l|l|}
		\hline
		&  &  &  &  &  \\ \hline
		&  &  &  &  &  \\ \hline
		&  &  &  &  &  \\ \hline
		&  &  &  &  &  \\ \hline
		&  &  &  &  &  \\ \hline
	\end{tabular}
\end{table}


\textbf{Modelo de Equação:}

\begin{equation} \label{eq:nome}
IC=\frac{F*9,81}{A}*10^{-6}
\end{equation}
em que,

IC - índice de cone,$MPa$;

F - força, $kgf$;

A - área do cone,$m^{2}$.

De acordo com a \cite{carvalho} 


Ilustrações e gráficos devem ser apresentados com tamanho e detalhes suficientes para a composição gráfica final, preferivelmente na mesma posição do texto.

\textbf{Gráficos:} devem apresentar-se sem bordas, descritos com o mesmo tipo e tamanho de letras contidas no texto e a legenda na posição inferior do mesmo. A numeração deve ser sucessiva em algarismos arábicos.

\textbf{Tabelas:} evitar tabelas extensas e dados supérfluos; adequar seus tamanhos ao espaço útil do papel e colocar, na medida do possível, apenas linhas contínuas horizontais; suas legendas devem ser concisas e autoexplicativas. Na discussão, confrontar os dados obtidos com a literatura. (Times New Roman, 11, Justificado).
\vspace{0.5cm}

\noindent\textbf{Modelos de Figuras:}


\begin{figure}[H]
\centering
\includegraphics[width=10cm,angle=0]{figures/grafico_1}
\caption{Mapas de teor de água das camadas de 0-2,2 e 0,2-0,4 m de profundidade.}
\label{fig:nome_referencia_figura1}
\end{figure}

\begin{figure}[H]
\centering
\includegraphics[width=10cm,angle=0]{figures/grafico_2.png}
\caption{Mapas do índice de cone (MPa) referente aos dados coletados nas diferentes profundidades nas linhas e nas entrelinhas da cultura da cana.}
\label{fig:nome_referencia_figura2}
\end{figure}


\noindent\textbf{Modelo de Tabela:}

\begin{table}[H]
\caption{Análise do IC nas linhas (L) e entrelinhas (E) de cana nas diferentes profundidades amostradas pelo índice de cone.}
\vspace{.2cm}
\centering
\resizebox{\textwidth}{!}{%
\begin{tabular}{ccccccccc}
\hline
Profundidade (m) & \multicolumn{2}{c}{0 a 0,1} & \multicolumn{2}{c}{0,1 a 0,2} & \multicolumn{2}{c}{0,2 a 0,3} & \multicolumn{2}{c}{0,3 a 0,4} \\ \hline
                 & L            & E            & L             & E             & L             & E             & L             & E             \\
Média (MPa)      & 1,39**       & 4,28**       & 1,86**        & 4,29**        & 2,20**        & 3,83**        & 2,46**        & 3,44**        \\
CV(\%)           & 54           & 57           & 55            & 54            & 46            & 49            & 48            & 43            \\ \hline
\end{tabular}%
}
\small{**:valores significativos para o nível de significância de 1\% pelo teste de Tukey; L – linhas; E – entrelinhas.}
\label{tab:nome_tabela1}
\end{table}

\section{CONCLUSÕES}

Neste capítulo devem ser apresentadas as principais conclusões ou considerações finais do trabalho.

Devem basear-se exclusivamente nos resultados do trabalho. Evitar a repetição dos resultados em listagem subsequente, buscando, sim, confrontar o que se obteve com os objetivos inicialmente estabelecidos.



%Referências Bibliográficas
%\textbf{REFERÊNCIAS}

As referências devem ser inseridas em ordem alfabética, alinhadas à esquerda, com fonte Calibri tamanho 10 e espaçamento simples. Usar Normas ABNT que podem ser acessadas no endereço: \url{https://www.normasabnt.org/referencias-bibliograficas/}

As referências devem ser listadas em ordem alfabética \cite{galvani}, \cite{national_instruments} e \cite{pandorfi}.




\bibliography{bibliografia.bib}
%\printbibliography[title={REFERÊNCIAS}]
%\printbibliography{bibliografia}


\section*{AGRADECIMENTOS}

A sessão de agradecimentos em um Resumo Expandido é o espaço onde você expressa sua gratidão às instituições, laboratórios e colaboradores que contribuíram de alguma forma para a realização da pesquisa. É imprescindível reconhecer as agências de fomento.

Inserir após as conclusões, de maneira sucinta. Se o projeto for financiado por alguma agência de fomento, citar a fonte. (Times New Roman, 11, Justificado).



\end{document}
